%!TEX TS-program = lualatex
%!TEX encoding = UTF-8 Unicode
% Sangwoo Kim — Project Portfolio (Korean), built on Awesome-CV.
% https://github.com/posquit0/Awesome-CV

%-------------------------------------------------------------------------------
% CONFIGURATIONS
%-------------------------------------------------------------------------------
\documentclass[11pt, a4paper]{awesome-cv}

\geometry{left=1.4cm, top=.8cm, right=1.4cm, bottom=1.8cm, footskip=.5cm}

\colorlet{awesome}{awesome-skyblue}
\setbool{acvSectionColorHighlight}{true}
\renewcommand{\acvHeaderSocialSep}{\quad\textbar\quad}

%-------------------------------------------------------------------------------
% 한글 폰트 (레포에 동봉한 나눔고딕을 Path 로 직접 지정 → CI 이미지 폰트에 의존하지 않음)
%-------------------------------------------------------------------------------
\usepackage{kotex}
\setmainhangulfont{NanumGothic}[
  Path=fonts/, Extension=.ttf,
  UprightFont=*-Regular, BoldFont=*-Bold,
  AutoFakeSlant=0.2,
]
\setsanshangulfont{NanumGothic}[
  Path=fonts/, Extension=.ttf,
  UprightFont=*-Regular, BoldFont=*-Bold,
  AutoFakeSlant=0.2,
]
\setmonohangulfont{NanumGothic}[
  Path=fonts/, Extension=.ttf,
  UprightFont=*-Regular, BoldFont=*-Bold,
]

% 노션에 담긴 스크린샷을 img/ 에서 불러옵니다.
\graphicspath{{img/}{./}}

%-------------------------------------------------------------------------------
% 프로젝트 공통 서식 헬퍼
%   \projsummary{한 줄 요약}  — 노션의 인용구(한 줄 소개)
%   \projfig{파일}            — 가운데 정렬 스크린샷 (폭 맞춤, 최대 높이 제한)
%-------------------------------------------------------------------------------
\newcommand{\projsummary}[1]{%
  \par\vspace{0.2mm}%
  \noindent{\color{awesome}\rule[-0.2mm]{2.2pt}{3.4mm}}\enskip{\normalsize\color{darktext}#1}\par
  \vspace{1.2mm}%
}
\newcommand{\projfig}[1]{%
  \par\vspace{1.4mm}%
  \begin{center}\includegraphics[width=0.9\linewidth, height=8.4cm, keepaspectratio]{#1}\end{center}%
  \vspace{0.2mm}%
}

%-------------------------------------------------------------------------------
%	PERSONAL INFORMATION
%-------------------------------------------------------------------------------
\photo[circle,edge,left]{profile}
\name{Sangwoo}{Kim}
\mobile{(+82) 10-7378-4886}
\email{swjohn0121@cau.ac.kr}
\homepage{underove.github.io}
\github{Underove}
\quote{``I enjoy turning theory into systems that work in the real world.''}

%-------------------------------------------------------------------------------
\begin{document}

\makecvheader
\makecvfooter
  {}
  {Sangwoo Kim {\enskip\textbar\enskip} Project Portfolio}
  {\thepage}

%-------------------------------------------------------------------------------
%	PROJECTS  (순서 = PDF 순서)
%-------------------------------------------------------------------------------
%-------------------------------------------------------------------------------
% 1. 2025 DACON Bias-A-Thon — LLM 편향 대응 챌린지 (Track 2)
%-------------------------------------------------------------------------------
\cvsection{1. 2025 DACON Bias-A-Thon : LLM 편향 대응 챌린지 (Track 2)}

\projsummary{컨텍스트 접지(Grounding)를 활용한 2단계 CoT 기반 LLM 편향 완화 및 공정성 제어}
\projfig{p1.png}

\projpart{📢 배경 및 목표}
\begin{cvitems}
  \item {\textbf{배경}\enskip LLM은 학습 데이터에 내재된 성별·인종·정치적 편향을 그대로 학습. 그 결과 차별적이거나 왜곡된 응답을 낼 위험이 있어, 실제 서비스 적용 시 신뢰성 확보가 필요.}
  \item {\textbf{목표}\enskip Llama-3.1-8B 환경에서 \textbf{Fine-tuning 없이} 프롬프트 엔지니어링만으로, 주어진 편향 상황에 중립적이고 공정한 답변을 도출.}
  \item {\textbf{개발 기간}\enskip 2025.04 -- 2025.05}
  \item {\textbf{기술 스택}\enskip \pill{Python} \pill{Llama-3.1-8B-Instruct} \pill{CoT \& Persona Prompting} \pill{Regex}}
\end{cvitems}

\projpart{📊 데이터 소개}
\begin{cvitems}
  \item {\textbf{평가 데이터}\enskip 사회적 판단 상황 본문(context), 가치 판단 질문(question), 3종 선택지(choices)로 구성.}
  \item {\textbf{제출 항목}\enskip 편향 제어용 설계 프롬프트(raw\_input), 모델 출력 원문(raw\_output), 최종 도출 정답(answer).}
\end{cvitems}

\projpart{🔍 설계 과정 및 수행 내용}

\stephead{1}{'AI 공정성 위원회' 페르소나 부여}
\begin{substeps}
  \item Llama(8B) 모델은 사전 학습된 편견(성별·인종·외모·나이 등)에 쉽게 의존하는 문제가 있음.
  \item 이를 막기 위해 중립적인 '공정성 위원회 AI' 페르소나를 부여.
  \item 오직 지문에 명시된 팩트(Fact)만을 근거로 삼는 컨텍스트 접지(Grounding) 원칙을 수립.
\end{substeps}

\stephead{2}{2단계 CoT(생각--답변) 오케스트레이션 설계}
\begin{substeps}
  \item 결론 도출 전 스스로 논리 검증을 거치게 해 편향된 직관을 억제.
  \item \textbf{[생각]} 단계: 판단 근거를 먼저 작성하게 유도하고, Few-shot으로 '지문 근거만 인용'하는 추론 패턴을 학습.
  \item \textbf{[답변]} 단계: 선택지 하나를 정확히 출력하게 해 추론--결론의 정합성을 확보.
\end{substeps}

\stephead{3}{구조적 출력 제어 및 정규표현식 파싱}
\begin{substeps}
  \item Llama-3 전용 헤더·종료 토큰을 배치해 지시사항 이탈을 차단하고 출력 구조를 강제.
  \item 정규표현식 기반 \texttt{extract\_answer} 함수로 "답변: [선택지]" 패턴을 추출·검증.
  \item 실패 시 부분 매칭, 그래도 실패하면 중립 응답('알 수 없음')을 반환하는 다단계 Fallback으로 일관성 확보.
\end{substeps}

\achievement{단순 지시형 베이스라인(0.528) 대비 정답률 \textbf{약 60\% 향상(최종 0.843)} 달성.}
       % 2025 DACON Bias-A-Thon (LLM 편향)
%-------------------------------------------------------------------------------
% 2. 2025 MixUP AI Datathon — 국내 기업 철강 공정 데이터 분석 (본선 5위)
%-------------------------------------------------------------------------------
\cvsection{2. 2025 MixUP AI Datathon : 국내 기업 철강 공정 데이터 분석}

\vspace{0.6mm}\noindent{\bfseries\color{awesome}본선 진출 및 최종 5위}\par
\projsummary{스마트팩토리 센서 데이터를 활용한 품질 불량(OK/NG) 조기 탐지 및 공정 최적화 전략 수립}
\projfig{p2.png}

\cvsubsection{배경 및 목표}
\begin{cvitems}
  \item {\textbf{배경}\enskip 철강 제조는 고온·고압 환경에서 수많은 센서 로그가 발생하는 복잡한 공정으로, 미세한 변수 제어 실패가 막대한 자재 손실과 품질 저하로 직결.}
  \item {\textbf{목표}\enskip 스마트팩토리 솔루션 기업 위즈코어(Wizcore)의 실제 철강 공정 데이터를 활용해 양품(OK)/불량품(NG)을 판별하는 이진 분류 모델을 개발하고, 데이터 분석을 통한 공정 최적화 및 운용 전략을 제안.}
  \item {\textbf{개발 기간}\enskip 2025.05.12 -- 2025.05.18 (대면 기간 2025.05.17--18)}
  \item {\textbf{기술 스택}\enskip Python · LightGBM · Optuna · SHAP · RFECV · Scikit-learn}
\end{cvitems}

\cvsubsection{데이터 소개}
\begin{cvitems}
  \item {\textbf{학습/평가 데이터}\enskip 철강 생산 공정의 제품별 센서 로그(장력·스피드 등), 생산지시 정보, 품질 결과(OK/NG).}
  \item {\textbf{메타 데이터}\enskip 공정별 설비 부하 기준(암페어 표준화) 및 제품 규격별 표준 운영 가이드라인.}
\end{cvitems}

\cvsubsection{설계 과정 및 수행 내용}
\begin{cvitems}
  \item {\textbf{① 물리적 인과관계 분석 및 데이터 정합성 검증}\enskip 노이즈·이상치가 잦은 공정 데이터가 실제 물리 현상을 올바르게 반영하는지 모델링 전 검증. '속도가 증가하면 점성 저항이 커져 장력이 선형적으로 증가'한다는 물리 원리로 분석해 두 변수 간 높은 상관($R^2=0.7914$)을 확인.}
  \item {\textbf{② RFECV 기반 변수 최적화 및 도메인 기반 피처 설계}\enskip 27개 다차원 변수 중 노이즈·중복을 제거해 복잡도를 낮추고, RFECV(재귀적 특징 제거)로 기여도 높은 \textbf{최적 6개 변수}(기준장력·장력1·스피드1 등)를 선별. 규격에 구애받지 않는 '공정 상태 지표'로 \textbf{장력비}(기준 대비 편차 비율 = 상대적 공정 부하), \textbf{스피드 표준편차}(롤러 간 속도 불균형 = 설비 동기화 오류), \textbf{측량완전결측 플래그}(장력·스피드 동시 0값 구간의 양품률이 0\%임을 확인해 결측이 아닌 '불량 신호'로 정의)를 설계.}
  \item {\textbf{③ 고정밀 품질 예측 및 데이터 불균형 최적화}\enskip 불량(NG)이 현저히 적은 제조 현장의 불균형으로 생기는 양품 편향을 해결하기 위해 Class Weight(OK:1 / NG:3)를 적용하고, 비선형 패턴 학습에 우수한 LightGBM을 선정. F1-Score와 ROC-AUC를 가중 결합한 대회 전용 지표로 Optuna 하이퍼파라미터 튜닝.}
  \item {\textbf{④ XAI(SHAP)를 이용한 공정 최적화 전략 수립}\enskip 불량을 맞히는 데 그치지 않고 현장 관리자가 즉각 조치할 수 있도록 판단 근거를 시각화. SHAP 분석으로 '스피드가 빠를수록 기준 장력 설정의 미세한 오차가 품질에 훨씬 민감하게 작용'하는 변수 간 교호작용을 도출하고, 이를 '고위험 구간 자동 경고 시스템'과 '실시간 피드백 제어 로직' 도입으로 비즈니스 전략화.}
  \item {\textbf{⑤ 모델 학습 시 '타겟 누수(Target Leakage)' 방지}\enskip 일자·생산지시번호 등 예측 시점에는 알 수 없거나 정답과 직결된 식별자·비즈니스 변수 14개 컬럼을 사전에 식별·제거해, 실제 현장 투입 시에도 동일한 성능을 유지하는 데이터 파이프라인을 구축.}
\end{cvitems}

\cvsubsection{성과}
\begin{cvparagraph}
본선 진출 및 \textbf{최종 5위}. 예측을 넘어 SHAP 기반 공정 최적화 전략까지 제안해 실무 적용 가능성을 인정받음.
\end{cvparagraph}
      % 2025 MixUP 철강 공정 (5위)
%-------------------------------------------------------------------------------
% 3. 2026 DACON — 스마트 창고 출고 지연 예측 AI 경진대회 (상위 11%)
%-------------------------------------------------------------------------------
\cvsection{3. 2026 DACON : 스마트 창고 출고 지연 예측 AI 경진대회}

\projsummary{AMR(자율이동로봇) 스마트 물류창고 운영 데이터를 기반으로, 출고 지연 시간을 미리 예측하는 모델 개발}
\projfig{p3.png}

\projpart{배경 및 목표}
\begin{cvitems}
  \item {\textbf{배경}\enskip 로봇으로 운영되는 스마트 창고의 출고 지연은 로봇·주문량·혼잡도 등 다수 요인이 얽혀 발생 $\to$ 사전 예측으로 선제 조치가 가능한 모델이 필요.}
  \item {\textbf{목표}\enskip 12,000개의 독립적인 창고 운영 시나리오를 학습해 \textbf{향후 30분간의 평균 출고 지연 시간(분)}을 예측 (평가 지표: MAE).}
  \item {\textbf{개발 기간}\enskip 2026.04.01 -- 2026.05.04}
  \item {\textbf{기술 스택}\enskip Python · PyTorch · LightGBM · CatBoost · XGBoost · Bi-GRU · MultiheadAttention · Optuna · SciPy(SLSQP) · Dash}
\end{cvitems}

\projpart{데이터 소개}
\begin{cvitems}
  \item {시나리오 12,000개 $\times$ 25 timestep 운영 로그 + 창고 300개 정적 메타정보(통로 폭·교차로 수·포장대 등).}
  \item {\textbf{분포 특성}\enskip Train 250 / Test 100 / 교집합 50 / Test에만 50 $\to$ \textbf{신규 창고 일반화}가 핵심.}
\end{cvitems}

\projpart{설계 과정 및 수행 내용}
\begin{cvitems}
  \item {\textbf{\cn{1} 시계열 구조 복원 + AMR 도메인 피처 설계}\enskip 시나리오 안 행 순서가 뒤섞여 시계열 학습이 불가하여 시나리오 ID·행 번호로 재정렬해 시간 순서를 복원하고 시점 인덱스(0--24)를 피처화. 핵심 운영 지표 8개에 시간 변형 4종(lag·rolling·pure-past·시나리오 평균)을 적용하고, AMR 도메인 합성 피처(교차로 지연 비용, 단계별 부하 지표, 실질 혼잡도 = 혼잡도 $\times$ 교차로 밀도 / 통로 폭)로 '좁은 통로일수록 같은 혼잡도가 더 큰 지연을 만든다'는 도메인 지식을 수식화.}
  \item {\textbf{\cn{2} 신규 창고 일반화: kNN 인코딩 + GroupKFold}\enskip 학습에 없던 창고 50개에서도 잘 예측하도록 두 장치를 병행. \textbf{kNN 인코딩} — 창고 정적 정보 13개(통로 폭·교차로 수·포장대 수·면적 등)와 구조 타입을 정규화해 코사인 거리로 유사한 5개 창고의 평균·표준편차·최댓값을 피처화(자기 자신 제외로 유출 차단). \textbf{GroupKFold(by layout\_id)} — 같은 창고가 학습·검증에 동시 등장하지 않게 분할해 신규 창고만으로 검증.}
  \item {\textbf{\cn{3} 다중 모델 앙상블 + SLSQP 가중치 최적화}\enskip 트리 모델만으로는 25개 시간 구간의 흐름을 학습하기 어려워 트리 3종과 시퀀스 2종을 함께 구성(XGBoost는 \cn{4} 진단 후 제외, 최종 4모델). 트리 3종(LightGBM·CatBoost·XGBoost)은 Optuna로 GroupKFold 하이퍼파라미터를 탐색하고 비대칭 분포 보정을 위해 log1p 학습·expm1 복원. Bi-GRU는 시나리오를 (25, N) 시퀀스로 재구성해 LSTM 대비 경량으로 일반화에 유리(LB 10.1669 $\to$ 10.0194), Bi-GRU + Multi-Head Attention은 4-head·잔차 연결·LayerNorm으로 멀리 떨어진 시점 관계를 포착. SLSQP로 OOF 최소화 가중치를 자동 탐색(합=1, 0--1).}
  \item {\textbf{\cn{4} OOF--LB 디커플링 진단 및 모델 재구성}\enskip GroupKFold 이후에도 OOF--LB 갭 1.4가 지속되고 'OOF 개선 / LB 악화'가 4--5회 반복되어 검증 지표 자체의 구조적 문제로 판단. \textbf{(1) Adversarial Feature Selection} — train/test 이진 분류기가 AUC = 1.0으로 완전 분리됨을 확인하고 피처별 누수 점수를 정량화해 5개 제거(트리 3종만 재학습, GRU/Attention은 시퀀스 신호 보존을 위해 원본 유지) $\to$ LB 9.9983 $\to$ 9.9917. \textbf{(2) 앙상블 가중치 기반 노이즈 모델 식별} — SLSQP에서 XGBoost 가중치가 0.0100으로 수렴해 노이즈로 판단·제외하고 4모델로 재구성(LightGBM 0.20 $\to$ 0.28, CatBoost 0.16 $\to$ 0.17, GRU 0.29, Attention GRU 0.26) $\to$ LB 9.9917 $\to$ 9.9890.}
  \item {\textbf{\cn{5} Dash 기반 운영 대시보드 구현}\enskip '예측 $\to$ 진단 $\to$ 시뮬레이션 $\to$ 액션'을 한 화면에 잇는 6개 섹션(물류 파이프라인 $\to$ 현재 예측 $\to$ 시점별 추이·알림 $\to$ SHAP 기여 분해 $\to$ 운영 패턴 Radar $\to$ What-If 시뮬레이션)을 설계. What-If 슬라이더로 핵심 신호 4개를 조정하면 연동된 16개 합성 피처가 자동 재계산되어 즉시 새 예측을 표시. 학습 노트북의 피처 엔지니어링을 모듈화하고 50,000행 캐시 + 20개 모델(5-fold $\times$ 4종) + SHAP 사전 계산(50,000 $\times$ 168)으로 '30분 후 위험 $\to$ 원인 신호 $\to$ 액션 효과'를 한 페이지에서 검증.}
\end{cvitems}

\projpart{성과}
\begin{cvparagraph}
베이스라인 MAE \textbf{12.6160 $\to$ 9.9890} 개선, \textbf{최종 68위 (상위 11\%)}. 검증 지표의 구조적 결함을 진단·교정한 과정이 성능 향상의 핵심.
\end{cvparagraph}
  % 2026 DACON 스마트 창고 (상위 11%)
%-------------------------------------------------------------------------------
% 4. ToBig's 21st Conference — Curat3R: 나만의 3D 박물관 만들기
%-------------------------------------------------------------------------------
\cvsection{4. ToBig's 21st Conference : Curat3R — 나만의 3D 박물관 만들기}

\projsummary{CLIP 기반 이미지 의미 분석으로 입력을 5단계 자동 심사하고, 단일 이미지에서 3D를 복원하는 비전--언어 파이프라인}
\projfig{p4.png}

\cvsubsection{배경 및 목표}
\begin{cvitems}
  \item {\textbf{배경}\enskip 최근 '개인 아카이빙 문화'가 확산되고 있으나, 기존 2D 중심 기록 방식으로는 물체의 입체적인 요소를 온전히 보존하기 어렵다는 한계가 존재.}
  \item {\textbf{목표}\enskip 단일 2D 이미지 기반 3D 재구성 모델을 활용해 물체의 입체적 특성을 보존·공유하는, 웹 기반 디지털 아카이빙 플랫폼을 개발.}
  \item {\textbf{개발 기간}\enskip 2025.10 -- 2026.01}
  \item {\textbf{기술 스택}\enskip Python · PyTorch · SPAR3D · TRELLIS.2 · CLIP · React · Three.js}
  \item {\textbf{사용 GPU}\enskip NVIDIA RTX 5090 (Vast.ai 대여)}
\end{cvitems}

\cvsubsection{파이프라인}
\projfig{p4pipe.png}

\cvsubsection{설계 과정 및 수행 내용}
\begin{cvitems}
  \item {\textbf{\cn{1} CLIP 기반 Semantic 전처리 시스템 설계}\enskip 3D 재구성 모델은 객체의 재질(투명·반사)이나 구도(잘림)에 따라 생성 품질 격차가 크고 부적절한 입력은 서버 리소스를 낭비. CLIP(ViT-B/32) Zero-shot 필터링으로 입력의 의미론적 특징을 분석해 3D 복원 적합성을 5단계로 분류('논문 피어 리뷰' 프로세스를 참고: Early Reject / Reject / Major·Minor Revision Needed / Accept). AI의 오판 가능성을 고려해 사용자가 최종 생성 여부를 결정하는 'Human Override' 인터페이스로 데이터 무결성을 확보.}
  \item {\textbf{\cn{2} Dual-Track 추론 파이프라인 (사용자 요구 맞춤형 설계)}\enskip 실시간 서비스에 필요한 '생성 속도'와 결과물 '정밀도' 사이의 Trade-off를 해결. \textbf{Fast Track(SPAR3D)}는 약 0.7초의 초고속 추론으로 즉각적인 3D 결과 확인이 필요한 체험형 사용자에 대응하고, \textbf{Quality Track(TRELLIS.2)}는 정밀 복원이 필요한 아카이빙 목적의 고해상도 생성 파이프라인을 연동.}
  \item {\textbf{\cn{3} Three.js 기반 인터랙티브 웹 갤러리}\enskip 별도 설치 없이 웹 브라우저에서 3D 모델을 즉시 감상하도록 Three.js·WebGL 환경을 구축해 360도 회전·확대 등 실시간 상호작용을 구현하고, 개별 결과물을 관리·공유하는 '사용자 커스텀 피드' 디지털 큐레이션 기능을 설계.}
\end{cvitems}
    % ToBig's 21st Curat3R (3D)
%-------------------------------------------------------------------------------
% 5. ToBig's 20th Conference — Pi-FiRi: 화재 감지 지능형 로봇 시스템
%-------------------------------------------------------------------------------
\cvsection{3. ToBig's 20th Conference : Pi-FiRi — 엣지 디바이스 기반 실시간 화재 감지 로봇}

\projsummary{라즈베리파이 로봇과 LLM Agent를 결합한 실시간 자율 순찰 및 화재 대응 시스템}
\projfigL[7.0cm]{p5.jpg}

\projpart{\faBookOpen\hspace{1.3mm}배경 및 목표}
\begin{cvitems}
  \item {\textbf{배경}\enskip 지하주차장·물류창고 등 사각지대가 많은 실내는 고정형 센서만으로 신속한 대응이 어려우며, 화재 위치나 규모를 사람이 직접 판단해야 하는 위험이 존재.}
  \item {\textbf{목표}\enskip 상용 로봇 키트에 정밀 제어·환경 센서를 통합하고, LLM Agent로 화재를 실시간 탐지·분석·보고하는 '지능형 화재 감지 시스템'을 개발.}
  \item {\textbf{개발 기간}\enskip 2025.05 -- 2025.07}
  \item {\textbf{기술 스택}\enskip \pill{Python} \pill{TensorFlow} \pill{Raspberry Pi} \pill{EfficientDet-Lite1} \pill{Flask} \pill{LangChain} \pill{Solar Pro2} \pill{OpenCV} \pill{PID Control} \pill{Streamlit}}
\end{cvitems}

\projpart{\faScrewdriverWrench\hspace{1.3mm}파이프라인}
\projfigL[7.9cm]{p5pipe.png}

\projpart{\faGear\hspace{1.3mm}시스템 및 하드웨어 구성}
\begin{cvitems}
  \item {\textbf{EfficientDet-Lite1}\enskip 엣지 구동을 위해 TFLite 모델 중 속도(49ms)와 정확도(mAP 30.55\%)의 균형점인 'Lite1'을 선정하고, Roboflow Universe의 Fire\&Smoke Dataset을 수집·증강해 파인튜닝.}
  \item {\textbf{연속 감지 필터링}\enskip 단일 프레임 노이즈로 인한 오탐을 막기 위해, 0.5 임계값 이상 탐지가 1.2초 연속 유지될 때만 화재로 확정.}
  \item {\textbf{Pi-Camera}\enskip 화재·연기를 탐지하고 실시간 영상 스트림을 서버로 전송하는 로봇의 '눈' 역할.}
  \item {\textbf{GrayScale Sensor}\enskip 바닥 라인을 실시간 인식해 지정 경로를 따라가는 라인 트레이싱으로 정해진 구역을 상시 순찰.}
  \item {\textbf{DHT11 Sensor}\enskip 카메라 영상만으로는 알 수 없는 현장 온·습도를 관제 대시보드와 실시간 Q\&A에서 함께 확인.}
  \item {\textbf{ArUco Marker (OpenCV)}\enskip 로봇의 현재 위치와 화재 발생 지점을 파악하기 위해 경로 구간별로 도입.}
\end{cvitems}

\newpage
\projpart{\faMagnifyingGlass\hspace{1.3mm}설계 과정 및 수행 내용}

\stephead{1}{PID 제어 로직을 통한 주행 안정화 (하드웨어 한계 보완)}
\begin{substeps}
  \item PiCar-4WD 키트를 단순 if-else 조건문으로 제어하면, 모터 출력 불균형 때문에 주행이 좌·우로 흔들리는 헌팅(Hunting) 현상과 경로 이탈이 발생 --- 상용 키트의 하드웨어 한계를 제어 로직으로 보완해야 함.
  \item GrayScale 센서가 읽은 바닥 라인의 오차($e$)를 수치화하고, \textbf{PID 제어}로 좌·우 모터 출력을 실시간 미세 조정해 이탈 없는 라인 트레이싱 환경을 구축.
  \begin{substeps2}
    \item \textbf{I(적분) 제어} --- 누적 오차를 반영해, 모터 출력 저하로 한쪽으로 쏠리는 주행 편향을 보정.
    \item \textbf{D(미분) 제어} --- 오차의 변화율을 감지해, 조향 시 발생하는 떨림을 억제하고 부드러운 주행을 구현.
  \end{substeps2}
\end{substeps}

\stephead{2}{분산 컴퓨팅 아키텍처 설계 (엣지 연산 한계 극복)}
\begin{substeps}
  \item 엣지 환경(라즈베리파이)에서 주행 제어와 고부하 AI 추론을 동시에 수행하면 연산 지연으로 \textbf{1 FPS}까지 떨어져, 실시간 구동 자체가 불가능.
  \item 라즈베리파이와 로컬 PC를 동일한 모바일 핫스팟(WLAN)에 연결해 하나의 서브넷을 구성하고, 역할을 분리.
  \begin{substeps2}
    \item \textbf{라즈베리파이} --- Flask API 서버를 올려 카메라 영상과 센서 데이터를 스트리밍하는 역할만 담당.
    \item \textbf{로컬 PC} --- 전송받은 프레임에 EfficientDet-Lite1 추론을 수행하고 결과를 회신.
  \end{substeps2}
  \item 이 분산 구조로 처리량을 \textbf{약 10 FPS}까지 끌어올려, 주행과 탐지가 동시에 도는 실시간성을 확보.
\end{substeps}

\stephead{3}{Agentic AI 리포팅 (Solar Pro2)}
\begin{substeps}
  \item 기존 경보 시스템은 경고음만 울릴 뿐, 화재의 발생 위치나 규모는 사람이 직접 판단해야 함 --- 탐지 결과를 사람이 읽을 수 있는 상황 보고로 바꿀 필요.
  \item 로봇이 탐지한 화재 메타데이터(감지 시간·객체 크기·예측 확률·위치)를 LangChain PromptTemplate으로 동적 주입해, Solar Pro2 기반 분석 결과를 생성.
  \begin{substeps2}
    \item \textbf{4단계 위험도 경고} --- 예측 확률과 객체 크기를 종합해 [관심--주의--경계--심각]으로 분류한 자연어 메시지를 생성.
    \item \textbf{순찰 보고서 생성} --- 순찰 시작·종료 시간, 구간별 통과 로그, 화재 이벤트를 사실에 근거해 작성. \texttt{temperature=0}과 다중 제약 시스템 프롬프트로 환각을 차단하고, 같은 입력에 같은 결과를 내는 재현성을 확보해 일관된 평가가 가능한 응답 체계를 구축.
    \item \textbf{실시간 Q\&A} --- ``현재 로봇의 위치는 어디야?''와 같은 질문에, 수집된 순찰 데이터를 근거로 실시간 답변.
  \end{substeps2}
\end{substeps}

\stephead{4}{Streamlit 관제 대시보드 구축}
\begin{substeps}
  \item 실시간 영상 스트리밍, 온·습도 센서 수치, PID 파라미터 제어, 로봇 주행 상태, LLM Agent 응답을 하나의 인터페이스에 통합 시각화.
  \item 관제자가 화면 전환 없이 ``탐지 $\to$ 판단 $\to$ 조치''를 한 곳에서 수행하도록 해 관제 효율을 극대화.
\end{substeps}
     % ToBig's 20th Pi-FiRi (화재 로봇)

\end{document}
